\documentclass[11pt,a4paper]{article}

% ------------------------------------------------------------------
% Conservation-Aware GNN Emulation of Stellar Nucleosynthesis
% Consolidated Research Landscape, Verified Evidence Base, and PhD Plan
% Synthesis of four research artifacts, June 2026.
% ------------------------------------------------------------------

\usepackage[margin=2.5cm]{geometry}
\usepackage{amsmath,amssymb}
\usepackage{microtype}
\usepackage{booktabs}
\usepackage{array}
\usepackage{tabularx}
\usepackage{pdflscape}
\usepackage{ragged2e}
\usepackage{enumitem}
\usepackage{xcolor}
\usepackage[colorlinks=true,linkcolor=blue!50!black,citecolor=blue!50!black,urlcolor=blue!50!black]{hyperref}

\newcolumntype{L}[1]{>{\RaggedRight\arraybackslash}p{#1}}

\newcommand{\Ye}{\ensuremath{Y_{\mathrm{e}}}}
\newcommand{\Msun}{\ensuremath{M_{\odot}}}
\newcommand{\code}[1]{\texttt{#1}}
\newcommand{\arxiv}[1]{\href{https://arxiv.org/abs/#1}{arXiv:#1}}

\title{\textbf{Conservation-Aware Graph Neural Network Emulation of Stellar Nucleosynthesis:}\\[0.4em]
\large Consolidated Research Landscape, Verified Evidence Base, and PhD Project Plan}
\author{Synthesis of four research artifacts\\
\small (landscape report, landscape audit, project plan, plan audit)}
\date{June 2026}

\begin{document}

\maketitle

\begin{center}
\begin{minipage}{0.9\textwidth}
\small\itshape
All audit corrections have been incorporated directly into the text; figures and citations below reflect the post-verification state of the evidence.
\end{minipage}
\end{center}

\begin{abstract}
\noindent
This report consolidates a full research cycle---landscape survey, adversarial audit of that survey, a detailed PhD project plan, and an adversarial audit of that plan---into a single document for a thesis on machine-learned emulation of stellar nucleosynthesis reaction networks. The headline conclusions are as follows. The bare concept ``graph neural networks for nuclear reaction networks'' is no longer novel: NuGNN (Kim, Chae, Ko, Mumpower \& Smith, \arxiv{2606.04491}, submitted 3~June 2026) is the first and so far only GNN surrogate for a nuclear-reaction-network \emph{solver}, demonstrated on a 690-isotope Type~I X-ray-burst network. However, the specific thesis niche remains open: a conservation-by-construction graph emulator for the silicon-burning / core-collapse-supernova-progenitor regime, benchmarked head-to-head against the published Grichener et al.\ 2025 fully-connected ``Nuclear Neural Network'' (NNN) on its public Zenodo data, and ultimately coupled in-the-loop into MESA via operator-split burning. No published or preprinted work combines hard baryon/charge conservation, a graph architecture, and an in-the-loop stellar-evolution demonstration for advanced burning.

The plan's evidence base survived two rounds of adversarial verification with no hallucinated citations and only minor factual corrections, all of which are integrated here. The audits did, however, surface one technical finding that materially reshapes the plan: the originally proposed flux-form architecture---predicting non-negative per-reaction fluxes and mapping them to abundance changes through the stoichiometric matrix---is exactly the catastrophic-cancellation operation that the quasi-statistical-equilibrium (QSE) literature has warned against in silicon burning since Hix \& Thielemann 1996. Conservation would be exact, but accuracy of the headline quantity (the electron fraction \Ye) would not follow. The revised plan therefore front-loads a 6--8 week QSE-cancellation kill-test, adopts signed-net-flux or projection-based conservation as the default architecture candidates, re-ranks the riskiest step from Phase~3 (MESA coupling, an engineering risk) to Phase~1 (the architectural bet, an existence risk), and converts the load-bearing collaboration with the Grichener/Renzo/Farmer group from assumed to contractual. The overall verdict across both audits: commit to the project, conditionally, with the de-risking spike and a written division of labor in place.
\end{abstract}

\tableofcontents

\section{Scientific Motivation: The Nucleosynthesis Bottleneck in Massive-Star Evolution}

A star above roughly $8\,\Msun$ ends its life in core collapse, and the outcome---neutron star versus black hole, explosion energy, nucleosynthetic yields---is largely set by the pre-collapse core structure. The controlling quantity is the electron fraction $\Ye \equiv \sum_i X_i Z_i / A_i$, because the effective Chandrasekhar mass scales as $M_{\mathrm{ch}}^{\mathrm{eff}} \propto \Ye^{2}$. During oxygen and silicon burning the core is reshaped by weak interactions---electron captures and $\beta$-decays---that are entirely absent from the small $\alpha$-chain networks used in most massive-star calculations. Farmer, Fields, Petermann et al.\ 2016 (ApJS 227:22, \arxiv{1611.01207}) quantified the stakes: at the onset of core collapse they find ${\approx}30\%$ variations in the central electron fraction and the mass locations of the main burning shells across network choices, and a minimum of ${\approx}127$ isotopes is needed to converge these values at the ${\approx}10\%$ level. Accurate \Ye, nuclear energy generation $\varepsilon_{\mathrm{nuc}}$, and pre-supernova neutrino losses require networks of hundreds of isotopes including neutron-rich species and tabulated weak rates (Suzuki et al.\ 2016).

The reason such networks are rarely used in full stellar-evolution runs is stiffness. The coupled abundance ODEs span reaction timescales differing by many orders of magnitude, forcing implicit integrators (backward Euler, Bulirsch--Stoer, semi-implicit midpoint) with repeated Jacobian construction and matrix inversions at every step. Coupled to the structure and mixing equations, cost and convergence risk explode, so late burning is routinely handled with ${\sim}21$-isotope hardwired $\alpha$-chain nets (the ``approx21 family,'' descended from Timmes et al.\ 2000) that get energy generation roughly right but \Ye{} and the neutrino physics wrong.

The state of the art and its mitigations define the playing field. MESA (Paxton et al.\ 2011--2019; Jermyn et al.\ 2023, ApJS 265:15) offers hardwired approximate nets and ``softwired'' nets in which all allowed pathways among the listed isotopes are linked (\code{mesa\_79}/\code{80}, \code{mesa\_127}, \code{mesa\_151}, \code{mesa\_160}, \code{mesa\_204}, the last tuned to reproduce the final \Ye{} of a 3298-isotope one-zone burn). Classical mitigation strategies include NSE/QSE group approximations (Hix \& Thielemann 1996), adaptive networks (Rauscher et al.\ 2002), operator-split burning that decouples the burn from the structure solve (\code{op\_split\_burn} in MESA, active at $T \gtrsim 3\times10^{9}$\,K, using a semi-implicit midpoint rule with adaptive sub-steps), and small-net evolution with large-net post-processing (the Sukhbold/Heger/Woosley KEPLER ecosystem, which should be treated as the dominant incumbent competitor to any emulator). The supporting code ecosystem comprises the \code{bbq} single-zone wrapper around MESA's network solver (Farmer 2023, Bulirsch--Stoer per Deuflhard 1983), the standalone networks SkyNet (Lippuner \& Roberts 2017), WinNet (Reichert et al.\ 2023), XNet and Torch, and \code{pynucastro} (Willcox \& Zingale 2018; Smith et al.\ 2023), all drawing on the JINA REACLIB rate database (Cyburt et al.\ 2010; the live JINA-CEE/CeNAM library now exceeds 60{,}000--80{,}000 rates depending on the library object queried). \code{pynucastro} is particularly important for this project: it exports the reaction network as a graph (NetworkX), provides REACLIB rate fits, tabulated weak rates, $Q$-values, partition functions, and separation energies---essentially the entire feature set and adjacency structure a graph emulator needs.

An emulator competing in this space must also be benchmarked against the \emph{non-neural} alternatives---NSE/QSE tabulation and lookup, Gaussian-process surrogates, and the KEPLER-style post-processing workflow---not only against other neural networks.

\section{The Two Anchor Works}

\subsection{Grichener et al.\ 2025 --- the published baseline (NNN)}

Grichener et al.\ 2025 (``Nuclear Neural Networks,'' ApJS 279, 49; \arxiv{2503.00115}; DOI 10.3847/1538-4365/ade717; authors Grichener, Renzo, Kerzendorf, Farmer, de Mink, Bellinger, Chan, Chen, Farag, Justham; lead author at Steward Observatory, University of Arizona) built a fully-connected MLP emulator of MESA r23.05.1's nuclear burning via \code{bbq}, for the softwired \code{mesa\_80} and \code{mesa\_151} networks. Inputs are $\log T$, $\log\rho$, and the composition vector; outputs are the post-step composition, a normalized nuclear-energy term, and a normalized neutrino-loss term. The architecture uses a first hidden layer of 1024 neurons and subsequent layers of 2048 (12 hidden layers for \code{mesa\_80}, 9 for \code{mesa\_151}), ReLU activations, L1 loss (log-scale composition, linear energy terms), Adam at $10^{-4}$, batch 512, and a log-softmax composition head---meaning the exponentiated outputs are normalized so that $\sum_i X_i = 1$; this enforces mass-fraction normalization only, not baryon-number or charge conservation under the reaction bookkeeping. A separate model is trained for each of nine logarithmically spaced timesteps from $10^{-6}$ to $100$\,s, with runtime interpolation between the two nearest. Training data: ${\sim}10^{6}$ samples per network, generated with \code{bbq} using a Sobol quasi-random sampler over $10^{9.2}\,\mathrm{K} < T < 10^{9.9}\,\mathrm{K}$ (${\approx}1.6$--$7.9$\,GK) and $10^{7} < \rho < 10^{9}\,\mathrm{g\,cm^{-3}}$ with $0.45 < \Ye < 0.5$, at a cost of a few${}\times10^{5}$ CPU-hours per regime. Critically, they found that sampling compositions only from real stellar models gave insufficient diversity, and broad---even unphysical---Sobol sampling was required for generalization.

Verified headline results: relative to the 22-isotope approx21 net, the NNN improves \Ye{} accuracy by 280--660\% (390--660\% for \code{mesa\_80}, 280--400\% for \code{mesa\_151}), average atomic/mass number by 150--360\%, and nuclear energy generation by 250--750\%. Neutrino losses improve by 100--$10^{6}$\% at short timesteps but become \emph{worse} than the small net for $\mathrm{d}t \gtrsim 0.1$\,s. The matched-context speedup over \code{bbq} is 20--34$\times$, and the NNN is even 1.5--4$\times$ faster than the 22-isotope net, with inference time roughly constant in $\mathrm{d}t$; these figures exclude the one-time training-data cost. The paper also discovered and fixed a real MESA bug---endo-energetic multi-particle rates miscomputed by 20--24 orders of magnitude from an omitted phase-space factor (their Appendix~B).

Acknowledged limitations, all verified: large per-isotope composition errors ($10^{-4} \lesssim \Delta X_i \lesssim 10^{-1}$, with worst-case predictions unrelated to the target); error accumulation over a 1000-step constant-$(T,\rho)$ rollout such that quantities other than \Ye{} become comparable to or worse than the small net; no extrapolation outside the training box (explicitly unsuitable for electron-capture or pulsational-pair-instability supernovae); energy terms emulated without thermodynamic consistency; no uncertainty quantification; and---decisive for this thesis---the MESA coupling is \emph{not implemented}; the work is an offline proof-of-concept validated against ${\sim}700$ held-out \code{bbq} test sets. The reproducibility package is public as Zenodo record 14873443 (MESA models, training/test sets, trained NNN models, analysis scripts); the complete 4\,TB of \code{bbq} runs (extra timesteps, extended density) is available on request only.

\subsection{NuGNN (Kim et al.\ 2026) --- the closest competitor}

NuGNN (Kim, Chae \& Ko, Sungkyunkwan University; Mumpower, Obsidian Research, Fort Wayne IN, and University of Notre Dame; Smith, Stellar Science Solutions, Lenoir City TN; \arxiv{2606.04491}, submitted 3~June 2026, v1 preprint) is the first GNN built directly on the bipartite isotope/reaction graph of a nuclear network: a surrogate for the PRISM solver on a 690-isotope network under general Type~I X-ray-burst conditions ($T \in [10^{6}, 2.5\times10^{9}]$\,K, $\rho \in [10^{-2}, 10^{8}]\,\mathrm{g\,cm^{-3}}$, $\Delta t \in [10^{-15}, 10^{-4}]$\,s---note the variable-$\Delta t$ input, in contrast to NNN's per-timestep models). The architecture uses heterogeneous isotope and reaction nodes, three connection types (isotope$\to$reaction, reaction$\to$isotope, isotope$\to$isotope), five message-passing blocks with graph attention, a seven-type reaction embedding, separation-energy node features, and a two-channel sigmoid output head. It was trained on 90{,}000 samples (72k/9k/9k split; LR $10^{-3}$ for 1{,}000 epochs then 200 at $10^{-4}$).

Accuracy: total signed-log MAE of 0.037 versus 0.11 (Res-U-Net) and 0.56 (fully-connected baseline)---a 3--10$\times$ edge---corresponding to ${\approx}1.8\%$ (all $X$) / 8.4\% (non-zero $X$) linear error on physically realistic data. Most importantly, when swapped into the network-evolution code in place of the solver, NuGNN reproduces the final abundance patterns (in-solver deviations mostly below ${\sim}15\%$) while both baseline architectures fail. This is the strongest existing evidence that the graph inductive bias pays off on nuclear networks.

Two properties define the open niche. First, NuGNN predicts $\Delta X$ directly and does \emph{not} enforce conservation by construction: it learns $\sum_i \Delta X_i \approx 0$ only implicitly from data, backed by a retry-with-smaller-$\Delta t$ heuristic, and the authors report trying and abandoning a conservation loss/output activation because their log-target preprocessing made it unstable. Second, its speed claim is hardware-asymmetric and must not be imitated: the original PRISM solver was timed on a CPU (Intel Xeon Platinum 8360Y) while the surrogates ran on an NVIDIA RTX 5090 GPU, with no clean speedup factor stated. Component inference times: 19.3\,ms (PyTorch) $\to$ 1.9\,ms (TensorRT) for NuGNN, versus 4.2\,ms for Res-U-Net (the slowest of the three, due to its heavy convolutional image-like inputs) and 0.1\,ms for the FNN (the fastest). NuGNN's advantage is accuracy, not raw per-call speed. Its total parameter count is not reported.

\section{The Wider Method Landscape}

\subsection{Other nucleosynthesis surrogates}

Three prior ML surrogates for astrophysical nuclear networks exist, all fully-connected/residual MLPs rather than GNNs: Fan et al.\ 2022 (\arxiv{2207.10628}; a 3-isotope carbon-burning network for Type~Ia flames in MAESTROeX), Grichener et al.\ 2025 (above), and Zhang/Yao et al.\ 2025 (DeePODE, \arxiv{2504.14180}, ApJ, DOI 10.3847/1538-4357/adf331)---DNN surrogates trained with evolutionary Monte Carlo sampling that replace stiff solvers in nuclear reacting-flow codes (Castro/MAESTROeX/FLASH context), tested on 3- and 13-isotope networks; its adaptive sampling, which concentrates training density where local gradients and timescales demand it, is the principled answer to where expensive \code{bbq} calls should be spent. On the r-process side, RHINE (Just, Xiong \& Mart\'inez-Pinedo, Phys.\ Rev.\ D 2026, \arxiv{2507.09040}; GSI/FAIR) is a neural emulator of r-process \emph{heating} coupled in-the-loop into neutron-star-merger hydrodynamics---but it tracks a handful of composition summary quantities rather than a full network, making it a complementary philosophy to contrast with rather than a competitor. Related but distinct: NN prediction of fission cross sections for r-process networks (Rodr\'iguez-S\'anchez et al.\ 2025, \arxiv{2502.17700}) and the ML nuclear-mass/rate literature.

Two adjacent GNN works sharpen the novelty statement. Padiyar, Dash \& Aditya (\arxiv{2603.22318}, March 2026) apply a GNN with species and reaction nodes and message-passing transformer layers to combustion \emph{mechanism reduction} (up to 95\% species/reaction reduction)---structurally the closest non-astro analogue to NuGNN, but not an abundance-evolution surrogate. A 2024 GNN learns nuclear cross sections across the chart of nuclides (\arxiv{2404.02332})---static property prediction, not network dynamics. The precise novelty claim is therefore: NuGNN is the only GNN \emph{surrogate for a nuclear-reaction-network solver}, and no work combines a conservation-enforcing GNN with an in-the-loop MESA demonstration for advanced burning.

\subsection{GNN physics simulators}

The methodological templates are Battaglia et al.\ 2018 (relational inductive biases), Sanchez-Gonzalez et al.\ 2020 (GNS, \arxiv{2002.09405}), and Pfaff et al.\ 2020 (MeshGraphNets, \arxiv{2010.03409}). Their transferable lessons: the encode--process--decode message-passing pattern; generalization across system size (GNS handles an order of magnitude more particles at test time than in training, because message passing is local and parameters are shared across nodes/edges); and the rollout-stability diagnosis---one-step-trained models never see their own accumulated error, so autoregressive rollouts drift off-distribution and diverge super-linearly, mitigated chiefly by corrupting training data with random-walk noise and by tuning message-passing depth. Dynami-CAL GraphNet (Sharma \& Fink, Nature Communications 17:1045, January 2026; \arxiv{2501.07373}) shows that enforcing pairwise conservation (linear and angular momentum) via edge-local frames yields stable long-rollout error and effective extrapolation. An important calibration from the audits: this evidence comes from \emph{non-stiff} particle/granular dynamics, and momentum is a different mathematical object from species mass; the transfer of ``noise injection plus conservation fixes rollout divergence'' to stiff networks near QSE is plausible but unproven, and should be treated as a hypothesis to test rather than an established remedy.

\subsection{Stiff-chemistry emulation: neural ODEs, operators, and conservation}

Combustion and astrochemistry have already solved several subproblems this thesis needs. ChemNODE (Owoyele \& Pal 2021) learns source terms as a neural ODE specifically to avoid integrate-time divergence of plain regressors. The autoencoder + latent-ODE family includes Grassi et al.\ 2022 (A\&A 668, A139), Sulzer \& Buck 2023 (\arxiv{2312.06015}; 29 species/224 reactions; 55$\times$ neural-ODE and up to 4000$\times$ linear-latent speedups), Phy-ChemNODE (with embedded elemental mass conservation), and---most directly relevant---MACE (Maes et al.\ 2024, ApJ 969:79, \arxiv{2405.03274}), which emulates an AGB-envelope network of 468 species and 6180 reactions and reports that enforcing element conservation as a \emph{soft} loss ``slowed down the training immensely'' because the decoder Jacobian must be computed---corroborating, together with NuGNN's abandoned conservation loss, that architectural (hard) conservation beats penalty terms.

Operator learning supplies the variable-$\Delta t$ mechanism that NNN's nine per-timestep models lack: DeepONet variants for stiff kinetics including a mass-conserving operator (Goswami et al.\ 2023, \arxiv{2302.12645}, CMAME 419:116674), extended FNO with element/mass constraints, AMORE (\arxiv{2510.12999}, mass-conserving adaptive multi-output operator), and Branca \& Pallottini 2024 (A\&A 684, A203) bringing DeepONet to ISM chemistry (9 species, 52 reactions; ${\approx}128\times$ speedup---$2\times10^{7}$ predictions in 25.36 CPU\,s versus 3240 CPU\,s for KROME). A crucial cautionary result: van de Bor et al.\ 2025 (\arxiv{2508.16114}) found the Branca \& Pallottini emulator prone to instability under iteration from accumulated prediction error---even operator-learned emulators are not rollout-stable for free; the known mitigations are noise injection (Holdship 2021) and training the iterative update directly (the MACE pattern).

The conservation-architecture lineage runs through the Chemical Reaction Neural Network (Ji \& Deng 2021, J.\ Phys.\ Chem.\ A 125:1082, \arxiv{2002.09062}), whose weights \emph{are} stoichiometric coefficients and Arrhenius parameters, and its atom-conserving extension (D\"oppel \& Votsmeier, Proceedings of the Combustion Institute 2024, 40(1), 105507), which embeds a conservation matrix---the reduced column echelon form of the null-space basis of the molecular matrix---as a fixed linear layer, improving training stability, speed, robustness to noise, and data efficiency. CRK-PINN and ANN-hard (Combustion and Flame 2024/2025, Xu group, SJTU) demonstrate that hard mass/energy/element constraints prevent the long-horizon error accumulation that plagues unconstrained surrogates. Supporting literatures to draw on: dedicated stiff-ML methods (Kim et al.\ 2021 stiff neural ODEs; Stiff-PINN; SPIN-ODE, \arxiv{2505.05625}), surrogate benchmarking and model selection (CODES, \arxiv{2410.20886}; systematic surrogate selection for nonequilibrium chemistry, \arxiv{2603.08567}), stellar-track emulation for reasoning about structure-equation feedback in the coupling phase (Hendriks \& Aerts 2019; Mombarg et al.\ 2021; Maltsev et al.\ 2024), equivariant GNNs and differentiable solvers, and the UQ/active-learning literature for fallback gating and training-set design.

\subsection{Novelty verdict}

GNNs applied to nuclear reaction networks are no longer novel as a bare concept but are minimally explored---one preprint, one regime. The open thesis space, confirmed by aggressive falsification attempts in both audits (no Grichener-group follow-up, no GNN-nucleosynthesis or MESA-coupling preprint exists as of mid-June 2026), is the \emph{combination} of: (i) the late-stage massive-star / CCSN-progenitor regime rather than X-ray bursts; (ii) hard baryon/charge (and energy-consistent) conservation built into a graph architecture; (iii) stable, validated in-the-loop MESA coupling over long timestep sequences; and (iv) variable-timestep generalization via neural-ODE/operator formulations on the graph. Every individual ingredient exists in some adjacent domain; the combination does not. The novelty window is actively closing---NuGNN appeared days before the landscape survey---so the framing must be re-checked against arXiv at each phase boundary.

\section{Technical Formulation}

\subsection{Graph design}

The natural representation is heterogeneous and bipartite: isotope nodes and reaction nodes, with reactant$\to$reaction and reaction$\to$product edges---the structure NuGNN validated and exactly what \code{pynucastro} exports. Because many nuclear reactions involve two or more reactants/products (triple-$\alpha$, $(\alpha,p)$ chains), reactions are formally hyperedges; the bipartite reaction-node construction is the standard tractable hypergraph encoding (cf.\ ChemHGNN, \arxiv{2506.11041}; Rxn-Hypergraph, \arxiv{2201.01196}) and is preferred here over explicit hypergraph convolutions.

Node features (isotopes): $Z$, $N$, $A$, mass excess/binding energy, one- and two-particle separation energies, temperature-dependent partition functions, current $X_i$ in log/asinh space, and learned $(Z,N)$ embeddings. Reaction-node/edge features: $Q$-value, REACLIB rate-fit coefficients or the evaluated log-rate at the current $T$ (and $\rho\Ye$ for tabulated weak rates), and a reaction-type embedding (REACLIB chapter / $(p,\gamma)$, $(\alpha,\gamma)$, $(\gamma,\alpha)$, $(p,\alpha)$, $(\alpha,p)$, $\beta^{+}$, e-capture). Global features: $T$, $\rho$, $\Delta t$. Output targets: composition change, nuclear energy generation, neutrino losses, with \Ye, $\bar{A}$, $\bar{Z}$ derived. Roughly four to six message-passing steps (NuGNN used five) is the starting depth.

\subsection{Conservation --- the central design decision, revised}
\label{sec:conservation}

Baryon number and charge are linear invariants of the abundance vector: every nuclear reaction $j$ satisfies $\sum_i A_i \nu_{ij} = 0$ and (with lepton bookkeeping for weak reactions) the appropriate charge balance, so any update of the form $\Delta Y = \nu F$ lies exactly on the conservation manifold regardless of the error in $F$. $\sum_i X_i = 1$ then follows from baryon conservation since $X_i = A_i Y_i$. This is the nuclear analogue of the D\"oppel--Votsmeier null-space layer and the decisive structural advantage over both NNN (normalization only) and NuGNN (soft, learned, abandoned). Two subtleties must be engineered in from the start. First, weak reactions change $Z$ at fixed $A$ and are precisely what drives the \Ye{} evolution the science cares about; the stoichiometric matrix must include lepton bookkeeping so that $A$ is conserved exactly \emph{while \Ye{} evolves physically}, and the conservation diagnostics must distinguish the two. Second---the critical audit finding---\emph{how} the fluxes are parameterized matters enormously in this regime:

\paragraph{The non-negative flux-and-subtract formulation is likely unsound for silicon burning.}
Above ${\sim}3$\,GK photodisintegration becomes important for all nuclei and the network sits in quasi-statistical equilibrium, where (per Hix \& Thielemann 1996/1998 and Guidry, Billings \& Hix 2011, \arxiv{1112.4738}) the net change in a species' abundance is orders of magnitude smaller than its formation or destruction rates, with fast equilibration reactions 6--8 orders of magnitude faster than the slow ones; MESA's own documentation flags the same near-cancellation hazard. The target band $10^{9.2}$--$10^{9.9}$\,K is mostly inside this regime. A network predicting forward and reverse fluxes to even 1--5\% accuracy therefore yields net $\Delta Y$ dominated by cancellation noise exactly where \Ye{} lives. The abundance-based competitors (NNN, NuGNN) sidestep this because the ground-truth solver already performed the subtraction. Conservation is preserved exactly either way---but conservation is not accuracy: the trajectory stays on the manifold while drifting to the wrong point on it.

The revised architecture candidates, in order:

\begin{enumerate}[leftmargin=2em]
  \item \textbf{Signed-net-flux GNN (primary).} Predict the \emph{signed net} per-reaction flux (forward minus reverse, as one quantity) and map through $\nu$. Conservation via the stoichiometric matrix still holds exactly; the forward/reverse cancellation never enters the network's output space. This preserves the conservation-by-construction headline while removing the conditioning pathology.
  \item \textbf{Net-$\Delta X$ with conservation projection (co-equal fallback).} Predict $\Delta X$ directly (the NuGNN-style target, which empirically works) and project onto the conservation manifold via a fixed null-space layer (the D\"oppel--Votsmeier construction applied to $(A, Z, \text{lepton})$ invariants). Slightly weaker physical interpretability, equally exact conservation.
  \item \textbf{GNN encoder + latent neural ODE} (the MACE pattern) for true continuous-$\Delta t$ evolution, retaining a conservation-respecting decoder---conservation in latent space is not automatic, so the projection/flux structure must live in the decode step.
  \item \textbf{GNN + DeepONet/operator head} with $\Delta t$ as the trunk input (Goswami 2023; Branca \& Pallottini 2024)---an alternative route to spanning $10^{-6}$--$10^{2}$\,s with one model, bearing in mind the documented iterative-instability risk.
  \item Graph attention / transformer layers---NuGNN already uses attention; an ingredient, not a differentiator.
\end{enumerate}

Which of (1) and (2) wins is an empirical question that the QSE kill-test (Section~\ref{sec:plan}, Phase~0/1) is designed to answer before any paper is promised. Hard constraints in the architecture: baryon number, charge/lepton bookkeeping, $\sum_i X_i = 1$. Soft constraints in the loss: energy consistency ($e_{\mathrm{nuc}}$ and neutrino terms should match $\sum_j Q_j F_j$ and the mass-excess accounting---predicting fluxes or net fluxes is what makes a physically consistent $e_{\mathrm{nuc}}$ possible at all), and optionally agreement with NSE/QSE equilibria at long $\Delta t$ and high $T$ as a physics-anchoring regularizer; the near-NSE structure of the Si-burning regime is a constraint that can \emph{help} the model as well as threaten it.

\subsection{Stiffness, rollouts, and timestep handling}

Remedies for autoregressive divergence, in priority order: training-noise injection (GNS random-walk noise---cheapest, highest expected impact, but unproven in the QSE regime); multi-step/pushforward rollout losses (backpropagating through $K$ predicted steps, equivalently training the iterative update directly); log- or asinh-space dynamics for abundances spanning 20+ decades (NuGNN's signed-log; D\"oppel--Votsmeier's asinh); and stiffness-aware time parameterization on a logarithmic $\Delta t$ grid pending the continuous-time head. Conservation removes one entire failure mode (drift of the invariants) but does not bound per-isotope, energy, or \Ye-evolution error---those are what the validation gates exist for.

\subsection{Size transfer}

Because message passing is local and parameters live on node/reaction-type embeddings rather than a fixed isotope index, a trained model can in principle be evaluated on a larger isotope set (\code{mesa\_80} $\to$ \code{mesa\_151} $\to$ \code{mesa\_204}) without retraining---the GNS/MeshGraphNets precedent. But that precedent concerns particle/mesh \emph{count}, not chemically new physics: the isotopes in \code{mesa\_151} $\setminus$ \code{mesa\_80} are neutron-rich species that are dynamically important precisely where \code{mesa\_80} fails. Expect zero-shot transfer to capture qualitative behavior and require fine-tuning on a modest \code{mesa\_151} set for per-isotope accuracy. This is a Phase-2 research deliverable and a hypothesis, not a freebie.

\section{Data and Compute Strategy}

\textbf{Year~1 lives entirely on public data.} Zenodo record 14873443 provides the \code{mesa\_80} and \code{mesa\_151} Sobol training/test grids at the nine fixed timesteps, the trained NNN models (the baseline to beat), the MESA stellar models defining the parameter space, and analysis scripts. MESA r23.05.1 and \code{bbq} are open; \code{pynucastro} supplies the graph structure, rates, weak-rate tables, and nuclear-data features.

\textbf{What must eventually be generated or obtained:} variable-$\Delta t$ training data (the public sets are at nine fixed timesteps); trajectory/off-trajectory rollout data for multi-step losses; new regimes (O-burning; lower-$T$ for electron-capture SNe; lower-$\rho$ for PPISN/PISN); and \code{mesa\_204}-scale sets. The 4\,TB of extra-timestep/extended-density \code{bbq} runs already exist on request from the Grichener group---negotiate access before generating anything.

\textbf{Compute arithmetic (verified by audit).} Grichener's few${}\times10^{5}$ CPU-hours per ${\sim}10^{6}$ samples implies ${\approx}0.1$--$0.5$ CPU-hr per sample. A 24-core personal machine running continuously yields ${\approx}17{,}000$ core-hours/month, i.e.\ ${\sim}35{,}000$--$170{,}000$ samples/month---an order of magnitude short of one fresh $10^{6}$-sample grid per month. Regenerating a Grichener-scale dataset for even one new regime is a multi-month, cluster-scale job requiring an institutional CPU allocation; this is the strongest single argument for living on public data in Year~1 and securing the 4\,TB set early. On the GPU side, a NuGNN-scale model (${\sim}10^{5}$--$10^{6}$ parameters, ${\sim}10^{5}$ training samples) trains in hours to one--two days on a single RTX Pro 6000 (96\,GB)---a reasoned estimate, not a cited benchmark---so the stated fleet of several RTX Pro 6000s plus an RTX 5070/5070~Ti is more than adequate for the offline phases, with hyperparameter sweeps rather than single runs dominating GPU time. Both the GPU fleet and the institutional CPU allocation are external dependencies that should be confirmed as actually available, not aspirational.

\textbf{Sampling design.} Adopt and extend Grichener's lesson: Sobol quasi-random sampling over the composition simplex under the \Ye{} constraint, deliberately including off-trajectory and unphysical states (the data-side counterpart of noise injection), and graduate to active/adaptive sampling---DeePODE's evolutionary Monte Carlo sampling is the published, nuclear-tested method for adapting sample density to local gradient and timescale structure. One validation caveat must be carried forward: error distributions measured on quasi-random compositions are not guaranteed to transfer to the thin, correlated manifold of real MESA trajectories (Grichener themselves restricted parts of their comparison to the approx21 isotope subset), so Sobol$\to$real-track transfer is an explicit validation target, not an assumption.

\section{Phased Plan (revised post-audit)}
\label{sec:plan}

Table~\ref{tab:phases} summarizes the phased plan with the audit-revised risk assignments.

\begin{landscape}
\begin{table}
\centering
\caption{Phased plan with revised risk assessment. Papers~1--2 form a standalone offline thesis spine if Phase~3 slips or is descoped to the hybrid result.}
\label{tab:phases}
\footnotesize
\begin{tabular}{@{}L{0.9cm}L{1.4cm}L{4.6cm}L{4.6cm}L{4.2cm}L{4.6cm}@{}}
\toprule
\textbf{Phase} & \textbf{Months} & \textbf{Goal} & \textbf{Deliverable} & \textbf{Riskiest assumption} & \textbf{Fallback} \\
\midrule
0 & 0--2 & Reproduce NNN baseline from Zenodo; build \code{mesa\_80} graph + stoichiometric/lepton matrix in \code{pynucastro}; \textbf{run the QSE-cancellation kill-test} & Benchmark harness; reproduced NNN metrics (within ${\sim}2\times$); go/no-go evidence on flux parameterizations at $\geq 3$--5\,GK & Public data/scripts are complete and runnable & Request missing pieces from Grichener group \\
\addlinespace
1 & 2--6 & Conservation-by-construction GNN (signed-net-flux or projection variant, per kill-test) on \code{mesa\_80}; offline head-to-head vs.\ NNN on identical data & \textbf{Paper 1}: machine-precision (${\sim}10^{-14}$ relative) $A$/$Z$ conservation + match-or-beat NNN per-isotope error distributions, Si-burning regime & The winning conservation variant reaches NNN-level \Ye{} accuracy in QSE conditions & If accuracy only ties NNN, conservation guarantee + regime + published-baseline comparison is still a publishable methods paper \\
\addlinespace
2 & 6--14 & Variable $\Delta t$ (neural-ODE/operator head), rollout training with noise injection and pushforward losses, \code{mesa\_80}$\to$\code{151} size transfer, UQ/OOD & \textbf{Paper 2}: variable-timestep, rollout-stable emulator with quantified transfer and uncertainty & Variable-$\Delta t$ data obtainable (4\,TB request) without a new \code{bbq} campaign & Train on 9 public timesteps + interpolation; defer continuous time \\
\addlinespace
3 & 12--24 & In-the-loop MESA coupling via \code{op\_split\_burn}; $20\,\Msun$ progenitor toward core collapse vs.\ large-net truth run & \textbf{Paper 3 (flagship)}: first GNN-emulated stellar evolution to core collapse with conservation guarantees & Numerical stability and energy-feedback control over $10^{3}$--$10^{4}$ in-loop steps & Hybrid fallback-to-solver run (emulator only in-distribution and gate-passing; \code{bbq} elsewhere)---still publishable \\
\addlinespace
4 & 24--36+ & One extension: ECSNe/PPISN regime, explosive burning, or defensive XRB head-to-head & \textbf{Paper 4} + thesis & Institutional CPU allocation for new-regime data & Astrochemistry cross-validation (MACE/UCLCHEM/CODES benchmarks, public data) as capstone \\
\bottomrule
\end{tabular}
\end{table}
\end{landscape}

\paragraph{The riskiest step, correctly identified, is Phase~1, not Phase~3.}
Phase~3 is an engineering risk and tractable: calling a serialized TorchScript model from MESA's Fortran is a solved problem in principle (FTorch, pytorch-fortran, TorchFort, all inference-only), and \code{op\_split\_burn} is the natural hook since it already evolves composition at fixed $T$, $\rho$ in adaptive sub-steps. The real engineering tail---float32/float64 casting at the ${\sim}10^{-14}$ conservation tolerance, batched per-cell inference across thousands of zones $\times\, 10^{3}$--$10^{4}$ steps, \code{op\_split\_burn}'s inability to supply Jacobian partials to the structure matrix solver---should be retired \emph{early and decoupled from the science run} by building the interop path against the $25\,\Msun$ split-burn test case that ships with MESA. Phase~1, by contrast, is an existence risk: if no conservation-respecting parameterization beats the MLP baseline on QSE-regime data, Papers~1--4 collapse together. Hence the kill-test in Phase~0 with an explicit threshold: \textbf{if both the signed-flux and net-$\Delta X$-with-projection variants fail to reach NNN-level \Ye{} accuracy at $\geq 5$\,GK, the conservation-first framing is not viable for Si-burning and the thesis pivots} to a lower-temperature regime or a different inductive bias.

\paragraph{Per-step error budget.}
Farmer 2016's ${\approx}10\%$ convergence figure is a model-to-model target, not an emulator tolerance, and propagating it over $10^{3}$--$10^{4}$ burn steps spans two orders of magnitude depending on the accumulation model: ${\sim}10^{-5}$/step under systematic (worst-case linear) accumulation up to ${\sim}10^{-3}$/step under a random-walk assumption. The plan adopts the conservative systematic-accumulation budget of ${\sim}10^{-5}$--$10^{-4}$ per-step \Ye{} error as the gate for attempting pure-emulator production runs, with the accumulation assumption stated explicitly. Energy error must additionally be small enough not to perturb structure-solver convergence---an empirical Phase-3 gate with its own dedicated test.

\paragraph{Timeline realism.}
Four papers in 36~months is optimistic once peer-review latency is included; the phase overlaps are designed so that Papers~1--2 form a standalone offline thesis spine if Phase~3 slips or is descoped to the hybrid result.

\section{Validation Methodology}

Validation proceeds on six fronts. (1) \emph{Rollout/error-accumulation tests}: reproduce Grichener's 1000-step constant-$(T,\rho)$ protocol first for apples-to-apples comparison---recognizing it as a baseline-parity test, not a stellar-trajectory test---then extend to varying-$(T,\rho)$ trajectories drawn from real MESA pre-supernova tracks, reporting per-step and cumulative error growth. (2) \emph{Conservation diagnostics}: baryon drift, charge drift, $\sum_i X_i - 1$, and energy consistency, with the flux/projection architecture demonstrating ${\sim}10^{-14}$ relative conservation while explicitly showing that \Ye{} still evolves correctly through weak reactions---the falsifiable headline claim NuGNN cannot make. (3) \emph{Ground truth}: \code{bbq}/MESA as primary, cross-checked against SkyNet and WinNet for solver-independent validation; report per-isotope error \emph{distributions}, not just means (the NNN's weak point), plus \Ye, $\bar{A}$, $\bar{Z}$, $e_{\mathrm{nuc}}$, and neutrino metrics at every timestep. (4) \emph{Speed}: report both matched-hardware and deployment-realistic timings, clearly labeled; contrast with Grichener's matched-context 20--34$\times$ and never quote a single GPU-vs-CPU ``speedup factor.'' (5) \emph{OOD/UQ}: test outside the training box, report graceful degradation, and implement an OOD flag (ensemble or MC-dropout on the flux head) feeding the fallback-to-solver gate. (6) \emph{Coupling-specific gates added by audit}: energy-generation error feeding back into the structure solve, behavior across the NSE transition, composition coupling in convective zones, and the Sobol$\to$real-track error transfer---none of which offline parity establishes on its own.

\section{Competitive Landscape, Differentiation, and Collaboration Risk}

\textbf{Active groups.} The Grichener/Renzo/Kerzendorf/Farmer/de Mink/Bellinger constellation (Arizona, MPA Garching, MSU, Technion-adjacent, and the MESA developer community) owns the NNN, the public dataset, \code{bbq}, and the CCSN-progenitor science case. The NuGNN team (Kim/Chae/Ko at Sungkyunkwan; Mumpower at Obsidian Research/Notre Dame; Smith at Stellar Science Solutions) owns the only graph surrogate and has demonstrated in-solver substitution for X-ray bursts. The GSI/FAIR group (Just/Xiong/Mart\'inez-Pinedo) owns the in-the-loop r-process heating niche with RHINE. The SJTU/Peking group (Zhang/Yao/Xu) drives DeePODE and the hard-conservation stiff-ODE methodology. Stony Brook (Zingale, \code{pynucastro}/AMReX) provides the tooling backbone; the astrochemistry emulation groups (Maes/MACE, Branca \& Pallottini, Holdship, Sulzer \& Buck, Grassi) are the methodological neighbors.

\textbf{Durable differentiation} rests on four legs: the Si-burning/CCSN-progenitor \emph{regime} (Grichener's turf, not NuGNN's); \emph{hard conservation} claimed first in Paper~1 (NuGNN tried and abandoned it); head-to-head comparison against a \emph{published} baseline on identical public data with matched-hardware timing; and \emph{MESA in-the-loop stellar evolution} as the flagship (versus NuGNN's one-zone network-evolution substitution). A reasoned forecast of the NuGNN group's moves---larger/r-process networks, improved conservation, new regimes possibly including CCSN---motivates shipping Paper~1 fast. The XRB head-to-head is defensive and optional; one cannot out-compete an entrenched group on their own benchmark.

\textbf{The collaboration is a single point of failure until formalized.} The plan's ``moat''---privileged access to \code{bbq}, the 4\,TB dataset, and MESA expertise---is also a hazard: the Grichener paper's own framing points toward integrating NNN-style emulators into stellar-evolution codes, i.e.\ toward the student's flagship Phase~3. The required mitigation is contractual, within the first three months: confirmed 4\,TB access, \code{bbq} pipeline access, and a written division of labor establishing the in-loop MESA coupling as the student's deliverable, with co-authorship agreed. If this cannot be formalized, Phase~3 is reclassified as high-risk and the offline Papers~1--2 spine becomes the primary thesis.

\textbf{Re-evaluation triggers.} If the NuGNN group posts a CCSN-regime or hard-conservation extension, accelerate Paper~1 and lean on the MESA-in-the-loop plus collaboration differentiation. If size transfer fails badly, pivot Paper~2 toward variable-$\Delta t$ + UQ and demote transfer to a fine-tuning result. If MESA coupling proves intractable within ${\sim}6$ months of Phase~3 start, freeze the hybrid fallback as Paper~3 and move to Phase~4. And at every phase boundary, re-run the novelty check against arXiv.

\section{Scope Extensions (ranked by payoff $\div$ difficulty)}

The core Si-burning hard-conservation GNN ranks first (zero new data, natural conservation structure, well-behaved near-NSE regime, no direct competitor). Then, in descending order: size transfer \code{mesa\_80}$\to$\code{151}$\to$\code{204} (public data plus a generated/negotiated \code{mesa\_204} set; the shared-embedding architecture enables it, new neutron-rich physics is the risk); variable-$\Delta t$ and in-the-loop MESA for the CCSN progenitor (the flagship---hardest, and most exposed to competition); electron-capture SNe (lower $T$, weak-rate-sensitive, \Ye-critical, sparse competition); PPISN/PISN (same framework, shifted parameter box); a defensive XRB head-to-head with NuGNN (their turf---only if forced); astrochemistry cross-validation on MACE/UCLCHEM-class problems with the CODES benchmarking framework (a low-risk methodological win that validates the conservation+rollout machinery on a published, well-understood problem before the thesis bet is settled); and r-process/kilonova networks (the graph scales, but RHINE occupies the in-the-loop heating niche with a complementary summary-variable philosophy worth contrasting rather than competing with).

\section{Condensed Reading List}

\textbf{Anchors:} Grichener et al.\ 2025 (ApJS 279, 49---read with the Zenodo package open) and Kim et al.\ 2026 (NuGNN, \arxiv{2606.04491}---architecture template and the critique of dense baselines). \textbf{Numerics and the QSE hazard:} Hix \& Thielemann 1996/1998 (Silicon Burning I/II); Guidry, Billings \& Hix 2011 (\arxiv{1112.4738}); Timmes 1999/2000. \textbf{GNN simulators:} Battaglia 2018; Sanchez-Gonzalez 2020 (GNS); Pfaff 2020 (MeshGraphNets); Sharma \& Fink 2026 (Dynami-CAL, Nat.\ Commun.\ 17:1045). \textbf{Conservation architectures:} Ji \& Deng 2021 (CRNN); D\"oppel \& Votsmeier 2024 (Proc.\ Combust.\ Inst.\ 40, 105507); CRK-PINN and ANN-hard (Combustion and Flame 2024/2025). \textbf{Stiff emulation and operators:} Owoyele \& Pal 2021 (ChemNODE); Goswami et al.\ 2023; Grassi et al.\ 2022; Sulzer \& Buck 2023; Maes et al.\ 2024 (MACE); Branca \& Pallottini 2024 with van de Bor et al.\ 2025 as its stability counterpoint; Zhang/Yao et al.\ 2025 (DeePODE). \textbf{In-the-loop exemplar:} Just, Xiong \& Mart\'inez-Pinedo 2026 (RHINE, PRD). \textbf{Physics and infrastructure:} Farmer et al.\ 2016; the MESA instrument papers culminating in Jermyn et al.\ 2023 (\code{op\_split\_burn}); \code{pynucastro} (Willcox \& Zingale 2018; Smith et al.\ 2023); Reichert et al.\ 2023 (WinNet); Lippuner \& Roberts 2017 (SkyNet); Cyburt et al.\ 2010 (REACLIB). \textbf{Benchmarking:} CODES (\arxiv{2410.20886}). \textbf{Adjacent GNN prior art to cite for precision:} Padiyar, Dash \& Aditya 2026 (\arxiv{2603.22318}) and the chart-of-nuclides cross-section GNN (\arxiv{2404.02332}).

\section{Caveats}

The field is moving fast: NuGNN appeared days before the first survey, it remains a v1 preprint with author-stated numbers, and a conservation-enforcing or CCSN-regime extension could appear at any time; the novelty framing must be revalidated continuously. Speedup figures across papers are not comparable (different hardware, network sizes, amortization of training cost), and NuGNN's in particular is component-timed and hardware-asymmetric. The QSE-cancellation severity for \emph{this specific architecture class} is an inference from well-established numerics, not a measured result---which is precisely why it is framed as a Phase-0 test rather than a verdict. Conservation bounds the invariants, not per-isotope, energy, or \Ye-evolution error. Size transfer is a hypothesis. The compute figures (0.1--0.5 CPU-hr/sample; one--two days per GPU training run; the $10^{-5}$--$10^{-4}$ per-step \Ye{} budget) are reasoned extrapolations, explicitly labeled. The reaction counts for \code{mesa\_80}/\code{mesa\_151} (the flux-head output dimension) should be confirmed via MESA's network introspection before sizing the architecture; softwired nets imply order several hundred to ${\sim}10^{3}$ reactions for \code{mesa\_151}---large but tractable. The 4\,TB dataset, the GPU fleet, and the institutional CPU allocation are all external dependencies, not facts. Finally, published versions of the most recent load-bearing sources (NuGNN, RHINE, Dynami-CAL) should be re-verified before any number is cited in the student's own papers.

\section{Conclusion}

Across five chats and four artifacts, the picture that emerges is consistent and, unusually for an AI-assisted research cycle, factually robust: two independent adversarial audits found no hallucinated citations in either the landscape survey or the project plan, and the corrections required were matters of completeness, framing, and a small number of details---all now folded into this document. The science case is genuine (network size controls \Ye{} and hence explodability, with a $\geq 127$-isotope convergence floor), the benchmark is public and reproducible (Grichener et al.\ 2025 plus Zenodo 14873443), the closest competitor (NuGNN) validates the graph inductive bias while leaving hard conservation, the Si-burning regime, and stellar-evolution coupling unclaimed, and the adjacent combustion/astrochemistry literature supplies tested solutions for conservation layers, variable timesteps, and rollout stability.

The decisive insight of the audit cycle is that the project's headline strength and its biggest technical risk were initially the same design choice. Conservation by construction is the right differentiator, but the non-negative flux-and-subtract realization of it walks directly into the QSE catastrophic-cancellation regime that defines silicon burning. The corrected plan keeps the differentiator and changes the realization: signed net fluxes or net-$\Delta X$ with a null-space conservation projection, selected by an explicit kill-test before Paper~1 is promised. With that test, a contractual division of labor with the Grichener group, early decoupled MESA-interop engineering, and the offline Papers~1--2 spine as a guaranteed fallback, the project is a sound, conditionally committed PhD: fast first paper on public data, a defensible and falsifiable conservation claim, a flagship in-the-loop demonstration no one has published, and clearly specified thresholds at which to pivot rather than persist.

\end{document}
